\section{Propuestas de Aplicación al Proyecto}

Esta sección detalla propuestas concretas para aplicar los conocimientos del marco teórico y literatura adicional al Portal Agro Comercial del Huila. Se estructuran por componentes clave del portal, indicando funciones apoyadas, problemas resueltos, mejoras recomendadas y oportunidades para el departamento.

\subsection{Integración de IoT y Agricultura de Precisión}

Basado en Elijah et al. (2018) y Pierce y Nowak (1999), proponemos integrar sensores IoT para monitoreo de suelos, clima y cultivos. Función: Dashboard de datos en tiempo real para productores.

Problema resuelto: Falta de información precisa, reduciendo intermediarios que controlan datos.

Mejoras: Usar redes LoRa para conectividad en áreas rurales del Huila, con algoritmos de IA para recomendaciones automáticas.

Recomendaciones: Capacitación en interpretación de datos; alianzas con proveedores de sensores.

Oportunidades: Aumento de productividad en fincas del Huila, atrayendo inversiones en tecnología agrícola.

Gráfica sugerida: Diagrama de flujo de datos IoT en el portal.

\subsection{Marketplace E-commerce y Marketing Digital}

Inspirado en Martínez y Sánchez (2018) y Gómez et al. (2021), el portal incluirá un marketplace con perfiles de productores, reseñas y pagos seguros.

Función: Conexión directa productor-consumidor, con herramientas de branding.

Problema: Intermediación en ventas, limitando precios justos.

Mejoras: Integrar storytelling digital y gamificación para engagement.

Recomendaciones: Certificaciones ecológicas integradas; soporte multilingüe para mercados internacionales.

Oportunidades: Posicionamiento del Huila como líder en agricultura digital, generando empleos en e-commerce.

Gráfica sugerida: Mapa de red de conexiones en el marketplace.

\subsection{Plataformas Tecnológicas y Trazabilidad}

De Ojeda-Beltrán (2022) y Mendoza (2024), proponemos una plataforma interoperable con APIs para integración de sistemas.

Función: Trazabilidad completa desde producción a venta.

Problema: Falta de confianza en calidad, perpetuando intermediarios.

Mejoras: Blockchain para registros inmutables; big data para predicciones de mercado.

Recomendaciones: Colaboraciones con entidades certificadoras; feedback loops para mejoras continuas.

Oportunidades: Acceso a mercados premium, aumentando exportaciones del Huila.

Gráfica sugerida: Arquitectura de plataforma con módulos de trazabilidad.

\subsection{Cierre de Brechas Digitales}

Según Torres (2019) y Ramírez (2021), implementar centro de capacitación virtual y soporte inclusivo.

Función: Cursos en línea y asistencia técnica.

Problema: Baja adopción por falta de habilidades.

Mejoras: Realidad aumentada para simulaciones; programas de inclusión para mujeres y jóvenes.

Recomendaciones: Alianzas con universidades; subsidios para dispositivos.

Oportunidades: Desarrollo humano en el Huila, reduciendo migración rural.

Gráfica sugerida: Gráfico de progreso en alfabetización digital.

\subsection{Innovación en Cadenas de Suministro}

De Wolfert et al. (2017) y Kamilaris et al. (2017), integrar optimización logística con IA.

Función: Matching inteligente de ofertas y demandas.

Problema: Ineficiencias logísticas favoreciendo intermediarios.

Mejoras: Integración con transportistas locales; predicciones de cosechas.

Recomendaciones: Pilotos en cadenas de valor del Huila; monitoreo de impacto.

Oportunidades: Eficiencia económica, liberando recursos para inversión.

Gráfica sugerida: Diagrama de cadena de suministro optimizada.

\subsection{Índice de Competitividad Digital}

Basado en UIT (2020) y CEPAL (2023), desarrollar dashboards de métricas para productores.

Función: Evaluación y mejora continua del índice digital.

Problema: Competitividad baja por brechas.

Mejoras: Benchmarks regionales; incentivos por mejoras.

Recomendaciones: Integración con políticas públicas; estudios de caso.

Oportunidades: Atracción de inversiones extranjeras al Huila.

Gráfica sugerida: Radar chart de indicadores de competitividad.

\subsection{Productos Ecológicos y Nichos de Mercado}

De López y García (2018) y Martínez (2021), sección dedicada a productos orgánicos.

Función: Plataforma para nichos premium.

Problema: Acceso limitado a mercados ecológicos.

Mejoras: Certificación digital automática; marketing enfocado.

Recomendaciones: Campañas de sensibilización; precios dinámicos.

Oportunidades: Valor agregado para agricultura sostenible del Huila.

Gráfica sugerida: Gráfico de crecimiento de mercado ecológico.